\documentclass[11pt, ukrainian]{article}
\setlength\textwidth{170mm} \setlength\hoffset{-25mm}%\setlength\textheight{277mm} \setlength\voffset{-38mm}
%\setlength\parindent{0mm}
\pagestyle{empty}
\usepackage[T2A]{fontenc}
\usepackage[utf8]{inputenc}
\usepackage[ukrainian]{babel}
\usepackage{graphicx}
\usepackage{caption}
\DeclareCaptionFormat{nocaption}{}
\captionsetup{format=nocaption,aboveskip=0pt,belowskip=0pt}
\usepackage[Export]{adjustbox}
\adjustboxset{max size={0.9\linewidth}{0.9\paperheight}}
\begin{document}
{04.11.2021 Модульна контрольна \No 1, варіант  \No \textbf { 60 }\par}
{
\begin {large}
{\begin {center}\textbf{Завдання 1}\end {center}}\par
Вказати значення виразу. \par
Дійсні значення (за наявності) округлити до 2 знаків після крапки. \par
За наявності помилок вказати \textbf{error}

\begin{center}
	\adjustimage{max size={0.9\linewidth}{0.9\paperheight}}
	{../pict/60_1.png}
\end{center}
\newpage

{\begin {center}\textbf{Завдання 2}\end {center}}\par
Вказати ТИП значення виразу з завдання 1. \par
За наявності помилок вказати \textbf{error}

\begin{center}
	\adjustimage{max size={0.9\linewidth}{0.9\paperheight}}
	{../pict/60_2.png}
\end{center}
\newpage

{\begin {center}\textbf{Завдання 3}\end {center}}\par
Вказати значення виразу. \par
Дійсні значення (за наявності) округлити до 2 знаків після крапки. \par
За наявності помилок вказати \textbf{error}

\begin{center}
	\adjustimage{max size={0.9\linewidth}{0.9\paperheight}}
	{../pict/60_3.png}
\end{center}
\newpage

{\begin {center}\textbf{Завдання 4}\end {center}}\par
Вказати ТИП значення виразу з завдання 3. \par
За наявності помилок вказати \textbf{error}

\begin{center}
	\adjustimage{max size={0.9\linewidth}{0.9\paperheight}}
	{../pict/60_4.png}
\end{center}
\newpage

{\begin {center}\textbf{Завдання 5}\end {center}}\par
Вказати значення виразу. \par
Дійсні значення (за наявності) округлити до 2 знаків після крапки. \par
За наявності помилок вказати \textbf{error}

\begin{center}
	\adjustimage{max size={0.9\linewidth}{0.9\paperheight}}
	{../pict/60_5.png}
\end{center}
\newpage

{\begin {center}\textbf{Завдання 6}\end {center}}\par
Вказати ТИП значення виразу з завдання 5. \par
За наявності помилок вказати \textbf{error}

\begin{center}
	\adjustimage{max size={0.9\linewidth}{0.9\paperheight}}
	{../pict/60_6.png}
\end{center}
\newpage

{\begin {center}\textbf{Завдання 7}\end {center}}\par
Вказати значення виразу. \par
Дійсні значення (за наявності) округлити до 2 знаків після крапки. \par
За наявності помилок вказати \textbf{error}

\begin{center}
	\adjustimage{max size={0.9\linewidth}{0.9\paperheight}}
	{../pict/60_7.png}
\end{center}
\newpage

{\begin {center}\textbf{Завдання 8}\end {center}}\par
Вказати ТИП значення виразу з завдання 7. \par
За наявності помилок вказати \textbf{error}

\begin{center}
	\adjustimage{max size={0.9\linewidth}{0.9\paperheight}}
	{../pict/60_8.png}
\end{center}
\newpage

{\begin {center}\textbf{Завдання 9}\end {center}}\par
Що буде виведено за виконання? \par
Повідомлення інтерпретатора про помилку позначати \textbf{error}

\begin{center}
	\adjustimage{max size={0.9\linewidth}{0.9\paperheight}}
	{../pict/60_9.png}
\end{center}
\newpage

{\begin {center}\textbf{Завдання 10}\end {center}}\par
Що буде виведено за виконання? \par
Повідомлення інтерпретатора про помилку позначати \textbf{error}

\begin{center}
	\adjustimage{max size={0.9\linewidth}{0.9\paperheight}}
	{../pict/60_10.png}
\end{center}
\newpage

{\begin {center}\textbf{Завдання 11}\end {center}}\par
Що буде виведено за виконання? \par
Повідомлення інтерпретатора про помилку позначати \textbf{error}

\begin{center}
	\adjustimage{max size={0.9\linewidth}{0.9\paperheight}}
	{../pict/60_11.png}
\end{center}
\newpage

{\begin {center}\textbf{Завдання 12}\end {center}}\par
Що буде виведено за виконання? \par
Повідомлення інтерпретатора про помилку позначати \textbf{error}

\begin{center}
	\adjustimage{max size={0.9\linewidth}{0.9\paperheight}}
	{../pict/60_12.png}
\end{center}
\newpage

{\begin {center}\textbf{Завдання 13}\end {center}}\par
Що буде виведено за виконання? \par
Повідомлення інтерпретатора про помилку позначати \textbf{error}

\begin{center}
	\adjustimage{max size={0.9\linewidth}{0.9\paperheight}}
	{../pict/60_13.png}
\end{center}
\newpage

{\begin {center}\textbf{Завдання 14}\end {center}}\par
Що буде виведено за виконання? \par
Повідомлення інтерпретатора про помилку позначати \textbf{error}

\begin{center}
	\adjustimage{max size={0.9\linewidth}{0.9\paperheight}}
	{../pict/60_14.png}
\end{center}
\newpage

{\begin {center}\textbf{Завдання 15}\end {center}}\par

Написати функцiю f, що отримує рядок i повертає ознаку того, що вiн складається рівно з однієї малої латинської літери. Стиль запису умови також оцінюється.\par
\newpage

\end {large}}
%end
\end{document}

